%!TEX root = ../../main.tex

\chapter{Gewichtsverfahren}
\label{chap:Gewichtsverfahren}

\section{EAN-13}
\label{sec:ean13}

Ein klassisches Beispiel für ein Gewichtsverfahren ist die Berechnung der Prüfziffer der \ac{EAN}, die insbesondere zur eindeutigen Identifikation von Handelsprodukten dient. Die EAN findet sich in der Regel sowohl im Klartext als auch in Form eines Strichcodes auf Produktverpackungen und Etiketten.

Die gebräuchliche EAN-13 besteht aus insgesamt 13 Ziffern, wobei die letzte Ziffer eine Prüfziffer darstellt. Diese wird mittels eines Modulo-10-Verfahrens ermittelt, um Eingabefehler (z. B. Zahlendreher oder Einzelfehler) zuverlässig zu erkennen.

Zur Berechnung der Prüfziffer werden die ersten 12 Ziffern positionsabhängig mit alternierenden Gewichten von 1 und 3 multipliziert. Anschließend werden die Produkte aufsummiert. Die Prüfziffer ergibt sich als diejenige Zahl, die zur gebildeten Summe addiert werden muss, um ein durch 10 teilbares Ergebnis zu erhalten. Formal ausgedrückt entspricht die Prüfziffer dem Wert, der notwendig ist, um die alternierend gewichtete Quersumme modulo 10 auf null zu ergänzen.

Die EAN-13 ist wie folgt aufgebaut:\newline
% EAN-13 Prüfzifferberechnung
\begin{equation}
v_0 v_1 v_2 v_3 v_4 v_5 v_6 v_7 v_8 v_9 v_{10} v_{11} v_{12}
\end{equation}

Für eine EAN-13 ist die Berechnung der Prüfziffer \( v_{12} \) wie folgt: \newline
\begin{equation}
\sum_{i \text{ gerade}} v_i + \sum_{i \text{ ungerade}} 3v_i \equiv 0 \pmod{10}
\end{equation}
Dieses Verfahren gewährleistet eine effiziente Plausibilitätskontrolle der EAN-Nummer im praktischen Einsatz \cite{Hoffmann2023-sv}.

\begin{figure}[H]
    \centering
    \begin{subfigure}{0.48\textwidth}
        \centering
        \includegraphics[width=0.6\linewidth]{ean_code.png}
        \caption{Beispiel eines EAN-13 Codes}
        \label{fig:ean13-example}
    \end{subfigure}
    \hfill
    \begin{subfigure}{0.48\textwidth}
        \centering
        \includegraphics[width=\linewidth]{ean_berechnung.png}
        \caption{Aufbau und Berechnung der EAN-13 Prüfziffer}
        \label{fig:ean13-structure}
    \end{subfigure}
    \caption{Darstellung und Aufbau des EAN-13 Code \cite{Hoffmann2023-sv}}
    \label{fig:ean13}
\end{figure}

\section{ISBN}

Ein ebenfalls weit verbreitetes Beispiel für ein Gewichtsverfahren ist die \ac{ISBN}, die in der Buchindustrie zur eindeutigen Identifikation von Büchern verwendet wird. 
Bis Ende 2006 wurde die ISBN-10 verwendet, die aus 10 Ziffern besteht und eine Prüfziffer enthält, die auf einem Modulo-11-Verfahren basiert.

Die ISBN-10 ist wie folgt aufgebaut: \newline
\begin{equation}
\label{eq:ISBN-10}
v_0 v_1 v_2 v_3 v_4 v_5 v_6 v_7 v_8 v_9,
\end{equation}

wobei die folgende Vorschrift für die Prüfziffer gilt: \newline
\begin{equation}
\label{eq:ISBN-10-Prüfziffer}
        10v_0 + 9v_1 + 8v_2 + 7v_3 + 6v_4 + 5v_5 + 4v_6 + 3v_7 + 2v_8 + 1v_9 \equiv 0 \pmod{11}
\end{equation}
Durch diesen Aufbau können Einzelfehler und Ziffernvertauschungen erkannt werden. Die mathematische Begründung hierfür lässt sich wie folgt darstellen: Durch die Verwendung der Primzahl 11 als Modul wird die Menge $\mathbb{Z}^*_{11}$ gebildet, die alle Zahlen zwischen 1 und 10 umfasst:
\begin{equation}
\mathbb{Z}^*_{11} = \mathbb{Z}_{11} \setminus \{0\} = \{1, 2, 3, 4, 5, 6, 7, 8, 9, 10\} 
\end{equation}
In dieser Menge sind ausnahmslos alle Gewichte enthalten, die in der Kongruenz (\ref{eq:ISBN-10-Prüfziffer}) verwendet werden. Ein grundlegendes mathematisches Prinzip besagt, dass ein Prüfziffercode mit der Beziehung $\sum_i g_i v_i \equiv 0 \pmod{p}$ genau dann in der Lage ist, alle Einzelfehler an einer Stelle $k$ zu erkennen, wenn das entsprechende Gewicht $g_k$ ein Element der Menge $\mathbb{Z}^*_p$ ist, also teilerfremd zum Modul $p$ ist.

Da bei der ISBN-10 alle Gewichte (1 bis 10) Elemente von $\mathbb{Z}^*_{11}$ sind und somit teilerfremd zu 11, erkennt der Code jeden Einzelfehler. Dies lässt sich mathematisch wie folgt begründen: Wenn eine Ziffer $v_i$ fehlerhaft als $v'_i$ gelesen wird, ändert sich die gewichtete Summe um $g_i \cdot (v'_i - v_i)$. Diese Änderung kann nur dann kongruent zu 0 modulo 11 sein, wenn entweder $g_i \equiv 0 \pmod{11}$ oder $v'_i - v_i \equiv 0 \pmod{11}$ gilt. Da $g_i$ teilerfremd zu 11 ist, bleibt nur $v'_i \equiv v_i \pmod{11}$ als Möglichkeit, was bedeutet, dass kein Fehler vorliegen kann.

Für die Erkennung von Vertauschungsfehlern gilt ein weiteres mathematisches Prinzip: Ein Prüfziffercode erkennt eine Vertauschung der Ziffern $v_k$ mit $v_l$ genau dann, wenn die Differenz der Gewichte $g_k - g_l$ ein Element der Menge $\mathbb{Z}^*_p$, also teilerfremd zum Modul $p$ ist. Bei der ISBN-10 ist jede Differenz zweier verschiedener Gewichte ein Element von $\mathbb{Z}^*_{11}$, da die Differenz zweier verschiedener Zahlen aus $\{1,2,\ldots,10\}$ nie durch 11 teilbar sein kann.

Mathematisch ausgedrückt: Bei einer Vertauschung zweier Ziffern $v_k$ und $v_l$ ändert sich die gewichtete Summe um $(g_k - g_l)(v_l - v_k)$. Da sowohl die Gewichtsdifferenz $g_k - g_l$ als auch die Zifferndifferenz $v_l - v_k$ (sofern die Ziffern unterschiedlich sind) nicht durch 11 teilbar sind, kann diese Änderung nicht kongruent zu 0 modulo 11 sein. Damit wird jede Vertauschung zweier unterschiedlicher Ziffern erkannt.

Ein besonderes Merkmal der ISBN-10 ist, dass als Prüfziffer auch der Wert 10 vorkommen kann, der dann durch das Zeichen 'X' dargestellt wird. Dies kann zu Verwirrung führen, da 'X' keine Dezimalziffer im eigentlichen Sinne ist und in manchen Systemen zu Verarbeitungsproblemen führen kann.

\begin{wrapfigure}{r}{0.4\textwidth}
    \centering
    \includegraphics[width=\linewidth]{isbn_aufbau.png}
    \caption{Beispiel für eine ISBN-13 \cite{Hoffmann2023-sv}}
    \label{fig:ISBN-13-Beispiel}
\end{wrapfigure}

Seit 2007 wird die ISBN-13 verwendet, die auf dem EAN-13-Code (vgl. Abschnitt \ref{sec:ean13}) basiert und dieselbe Struktur mit alternierenden Gewichten (1 und 3) und Modulo 10 verwendet. Ein Beispiel für eine ISBN-13 mit den Bedeutung der einzelnen Ziffern ist in Abbildung \ref{fig:ISBN-13-Beispiel} dargestellt.

Die ISBN-13 verwendet die Präfixe 978 oder 979-1 bis 979-9 und integriert dadurch die ISBN vollständig ins EAN-System. Die Umstellung von ISBN-10 auf ISBN-13 erfolgte primär aus praktischen Erwägungen: Durch die Erweiterung des Nummernraums können deutlich mehr Bücher eindeutig identifiziert werden. Zudem ermöglicht die Integration in das EAN-System eine nahtlose Verarbeitung in bestehenden logistischen Prozessen des Handels. Aus mathematischer Sicht ist jedoch die ISBN-10 mit ihrem Modulo-11-Verfahren dem neuen Format überlegen. Während die ISBN-10 durch ihre Primzahlbasis alle Vertauschungsfehler erkennt, kann die ISBN-13 mit ihrem Modulo-10-Verfahren bestimmte Ziffernvertauschungen nicht identifizieren. Die Vereinheitlichung mit dem EAN-System hat somit zu einer Verschlechterung der fehlererkennenden Eigenschaften geführt – ein klassischer Fall, bei dem praktische Anforderungen und theoretische Optimalität gegeneinander abgewogen werden mussten \cite{Hoffmann2023-sv}.

\section{Luhn-Algorithmus}

Der Luhn-Algorithmus oder auch Modulo-10-Algorithmus ist ein weiteres Beispiel für ein Gewichtsverfahren und wird verwendet, um eine Prüfzimmer zu berechnen, mit der sich unter anderem Kredidkartennummern,Indentifikationsnummern oder andere finanzielle Kennnummern valieren lassen.
Er wurde bereits 1954 von dem deutschen Informatiker Hans Peter Luhn patentiert und seitdem angewendet.

Im Folgenden wird der Luhn-Algorithmus exemplarisch anhand von Kreditkartennummern erläutert, da diese zu den verbreitetsten Anwendungsfällen zählen. Eine Kreditkartennummer besteht aus einer vierzehnstelligen Basisnummer sowie einer zu berechnenden Prüfziffer~$p_0$. Die Ziffern seien durch $v_i \in \{0, \dots, 9\}$ bezeichnet, wobei $i = 0, \dots, 13$. \newline
\begin{equation}
\label{eq:Luhn}
v_0 v_1 v_2 v_3 v_4 v_5 v_6 v_7 v_8 v_9 v_{10} v_{11} v_{12} v_{13} \, p_0
\end{equation}
Zur Berechnung der Prüfziffer~$p_0$ wird folgendermaßen vorgegangen:

\begin{enumerate}
  \item Zähle die Stellen der Basisnummer von rechts nach links mit den Indizes $i = 13, 12, \dots, 0$.
  \item Jede zweite Ziffer beginnend von rechts (also mit ungeradem Index $i$ im ursprünglichen Indexsystem) wird verdoppelt.
  \item Ist das Ergebnis der Verdopplung zweistellig, so wird die Quersumme dieser Zahl gebildet, d.\,h. die Summe der Ziffern von $2v_i$.
  \item Alle bearbeiteten und unbearbeiteten Ziffern werden zur Gesamtsumme~$S$ addiert.
  \item Die Prüfziffer~$p_0$ wird so gewählt, dass $S + p_0$ ein Vielfaches von~$10$ ist.
\end{enumerate}

Die Gesamtsumme~$S$ ergibt sich damit formal zu: \newline \begin{equation}
\label{eq:Luhn-Summe-fallunterscheidung}
S = \sum_{i=0}^{13}
\begin{cases}
v_i, & \text{wenn $i$ gerade (von rechts gezählt)} \\
2v_i, & \text{wenn $i$ ungerade und } 2v_i < 10 \\
\displaystyle \sum_{j=0}^{1} z_j, & \text{wenn $i$ ungerade und } 2v_i \geq 10;\ z_j \text{ als Ziffern von } 2v_i
\end{cases}
\end{equation}

Alternativ lässt sich die Gesamtsumme auch in kompakterer Notation durch eine verschachtelte Summe darstellen: \newline
\begin{equation}
\label{eq:Luhn-Summe-quersumme}
S = \sum_{\substack{0 \leq i \leq 13 \\ i\ \text{gerade (von rechts)}}} v_i
\ + \ 
\sum_{\substack{0 \leq i \leq 13 \\ i\ \text{ungerade (von rechts)}}}
\left(
\begin{cases}
2v_i, & \text{falls } 2v_i < 10 \\
\displaystyle \sum_{j=0}^{1} z_j, & \text{sonst}
\end{cases}
\right)
\quad \text{mit } z_j \text{ als Ziffern von } 2v_i
\end{equation}
Die Prüfziffer~$p_0$ ergibt sich abschließend durch: \newline \begin{equation}
\label{eq:Luhn-Pruefziffer}
p_0 = (10 - (S \bmod 10)) \bmod 10
\end{equation}
Gegeben sei die 14-stellige Kreditkartennummer: \newline \begin{equation*}
v = 4\ 9\ 1\ 6\ 4\ 9\ 4\ 0\ 3\ 1\ 4\ 1\ 7\ 2
\end{equation*}
Die Ziffern werden von rechts nach links mit den Indizes $i = 13$ (linke Ziffer) bis $i = 0$ (rechte Ziffer) nummeriert:

\begin{center}
\begin{tabular}{|c|c|c|c|c|}
\hline
$i$ & $v_i$ & Position (von rechts) & $2v_i$ & Anzuwendender Wert \\
\hline
13 & 4 & 14 (gerade) & –  & 4 \\
12 & 9 & 13 (ungerade) & 18 & 9 (1+8) \\
11 & 1 & 12 (gerade) & –  & 1 \\
10 & 6 & 11 (ungerade) & 12 & 3 (1+2) \\
9  & 4 & 10 (gerade) & –  & 4 \\
8  & 9 & 9 (ungerade)  & 18 & 9 (1+8) \\
7  & 4 & 8 (gerade) & –  & 4 \\
6  & 0 & 7 (ungerade)  & 0  & 0 \\
5  & 3 & 6 (gerade) & –  & 3 \\
4  & 1 & 5 (ungerade)  & 2  & 2 \\
3  & 4 & 4 (gerade) & –  & 4 \\
2  & 1 & 3 (ungerade)  & 2  & 2 \\
1  & 7 & 2 (gerade) & –  & 7 \\
0  & 2 & 1 (ungerade)  & 4  & 4 \\
\hline
\end{tabular}
\end{center}

Die Summe der anzuwendenden Werte ergibt: \newline \begin{equation*}
S = 4 + 9 + 1 + 3 + 4 + 9 + 4 + 0 + 3 + 2 + 4 + 2 + 7 + 4 = 56
\end{equation*}
Die Prüfziffer~$p_0$ berechnet sich zu: \newline \begin{equation*}
p_0 = (10 - (56 \bmod 10)) \bmod 10 = (10 - 6) \bmod 10 = 4
\end{equation*}
Die vollständige Kreditkartennummer lautet somit: 4916\ 4940\ 3141\ 724.

\todo{hier noch schreiben, warum das gut ist ? Erkennt man damit gut fehler oder wie ist das zu verstehen?}